(1) 线性浅水连续方程 \(\eta_t+h u_x=0\)、动量方程 \(u_t+g\eta_x=0\) 给 \(\eta_{tt}=gh\,\eta_{xx}\)。故两侧波速 \(c_i=\sqrt{gh_i}\)，色散关系 \(\omega=c_i|k|\)，相速和群速的大小在非色散浅水极限都等于 \(c_i\)。
(2) 右行波的速度扰动满足 \(hu=c\eta\)，左行满足 \(hu=-c\eta\)。在 \(x=0\) 用水位连续 \(\eta_i+\eta_r=\eta_t\) 与流量连续 \(c_1(\eta_i-\eta_r)=c_2\eta_t\)，得到 \(r=(c_1-c_2)/(c_1+c_2)\)、\(t=2c_1/(c_1+c_2)\)，均为水位振幅系数。
(3) 浅水线性波平均能流为 \(\overline P=\frac12\rho g c\,|\eta_0|^2\)（每单位横向宽度），所以反射能流比 \(R=|r|^2\)、透射能流比 \(T=(c_2/c_1)|t|^2=4c_1c_2/(c_1+c_2)^2\)。代数恒等式 \((c_1-c_2)^2+4c_1c_2=(c_1+c_2)^2\) 给 \(R+T=1\)。
(4) 若 \(h_2=h_1\)，则 \(r=0,t=1\)；若 \(h_2\to0^+\)，则 \(r\to1,t\to2\) 但透射能流趋零；若 \(h_2\to\infty\)，则 \(r\to-1,t\to0\)。后两极限通常使某侧不再满足相同浅水条件，故仅用作数学一致性检验，不能据此断言真实极深水体仍传播浅水波。
